\chapter{PENDAHULUAN}
\section{Latar Belakang Masalah}


Tuliskan di sini \textbf{Pedoman Penulisan:} Latar belakang masalah berisi garis besar alasan pembuatan skripsi dan permasalahan yang ada sehingga membutuhkan aplikasi dalam penyelesaian permasalahan.



\section{Identifikasi Permasalahan}


Tuliskan di sini \textbf{Pedoman Penulisan:} Berisi garis besar alasan pembuatan skripsi dan permasalahan yang ada sehingga membutuhkan aplikasi dalam penyelesaian permasalahan.



\section{Perumusan Masalah}


Tuliskan di sini \textbf{Pedoman Penulisan:} Perumusan masalah berbentuk kalimat tanya berdasarkan masalah yang akan dibahas.



\section{Maksud dan Tujuan}


Tuliskan di sini \textbf{Pedoman Penulisan:} Menjelaskan tujuan membuat Skripsi ini apa? Serta manfaat yang di berikan dari penulisan skripsi ini.\\
\textbf{Contoh Tujuan:}\\
1. Membangun aplikasi berbasis Android yang mampu membantu mengefektifkan waktu tunggu customer di bengkel mobil.\\
\textbf{Contoh Manfaat:}\\
1. Manfaat untuk penulis: Sebagai salah satu syarat kelulusan Program Sarjana pada Program Studi Teknologi Informasi UBSI.\\
2. Manfaat untuk objek penelitian: Sebagai alternatif peningkatan mutu pelayanan terhadap customer.\\
3. Manfaat untuk pembaca: Memberikan pemahaman mengenai pembuatan aplikasi antrian.



\section{Metode Penelitian}
\subsection{Teknik Pengumpulan Data}
\subsubsection*{A. Observasi}


Tuliskan di sini \textbf{Pedoman Penulisan:} Menjelaskan tentang proses observasi yang dilakukan.


\subsubsection*{B. Wawancara}


Tuliskan di sini \textbf{Pedoman Penulisan:} Menjelaskan proses wawancara yang dilakukan.


\subsubsection*{C. Studi Pustaka}


Tuliskan di sini \textbf{Pedoman Penulisan:} Menjelaskan tentang studi pustaka dan sumber referensi yang dipakai.



\subsection{Model Pengembangan Software}


Tuliskan di sini \textbf{Pedoman Penulisan:} Model Pengembangan software disesuaikan dengan keinginan yang dipilih oleh penulis, seperti model Waterfall, Spiral, Incremental, Prototype, Rapid Application Development (RAD), dll.


\subsubsection*{A. Analisa Kebutuhan Software}
\subsubsection*{B. Desain}
\subsubsection*{C. Code Generation}
\subsubsection*{D. Testing}
\subsubsection*{E. Support}

\section{Ruang Lingkup}


Tuliskan di sini \textbf{Pedoman Penulisan:} Penulis menjelaskan proses-proses yang dibahas pada saat pembuatan dan pengembangan aplikasi berdasarkan latar belakang permasalahan.


